\section{Instalación de Docker (Ubuntu 24.04)}
\begin{itemize}
    \item[§] \textbf{Desinstalación de versiones antiguas}\\
Para desinstalar versiones antiguas de docker 
\begin{lstlisting}
for pkg in docker.io docker-doc docker-compose docker-compose-v2 podman-docker containerd runc; do sudo apt-get remove $pkg; done
\end{lstlisting}
en este caso se vería la desinstalación de los paquetes típicos de docker.
    \item[§] \textbf{Añadir respositorio Docker para Ubuntu}
\begin{lstlisting}
# Para añadir la clave GPG oficial de Docker
sudo apt update
sudo apt install -y ca-certificates curl
sudo install -m 0755 -d /etc/apt/keyrings
sudo curl -fsSL https://download.docker.com/linux/ubuntu/gpg -o /etc/apt/keyrings/docker.asc
sudo chmod a+r /etc/apt/keyrings/docker.asc

# Añadiendo el repositorio a las fuentes de APT
echo \
  "deb [arch=$(dpkg --print-architecture) signed-by=/etc/apt/keyrings/docker.asc] https://download.docker.com/linux/ubuntu \
  $(. /etc/os-release && echo "$VERSION_CODENAME") stable" | \
  sudo tee /etc/apt/sources.list.d/docker.list > /dev/null
sudo apt update
\end{lstlisting}
La clave GPG de Docker se utiliza para firmar los paquetes y repositorios, garantizando de esta manera la integridad y autenticidad de las imágenes y software descargados.\\\\Las fuentes de APT son los repositorios, servidores remotos o locales, desde los que el gestor de paquetes APT descarga e instala software. Suelen definirse en 
    \begin{itemize}
        \item \textit{/etc/apt/sources.list-archivo principal}
        \item \textit{/etc/apt/sources.list.d/--directorio para fuentes adicionales (archivos .list o .sources)}
    \end{itemize}
    cada una de las líneas tiene el formato \textit{deb http://URL/ distribución componentes}
    \begin{table}[H]
        \begin{center}
            \begin{tabular}{|l||l|}\hline
                \textbf{Campo\hs{30}} & \textbf{Significado\hs{200}} \\\hline\hline
                deb & Paquetes binarios precompilados \\\hline
                deb-src & Código fuente \\\hline
                http://URL/ & Ubicación del repositorio (también \textit{file://}, \textit{ftp://}, \textit{cdrom:}) \\\hline
                distribución & Rama o versión (\textit{noble}, \textit{stable}, etc) \\\hline
                componentes & \textit{main}, \textit{restricted}, \textit{universe}, \textit{multiverse} \\\hline
            \end{tabular}
        \end{center}
    \end{table}
    \item[§] \textbf{Instalación de los paquetes de Docker Engine}
\begin{lstlisting}
sudo apt install -y docker-ce docker-ce-cli containerd.io \
docker-buildx-plugin docker-compose-plugin
\end{lstlisting}
    \item[§] \textbf{Asignación de permisos al usuario para ejecutar docker cli}
\begin{lstlisting}
sudo groupadd docker
sudo usermod -aG docker $USER
newgrp docker
\end{lstlisting}
    \item[§] \textbf{Configuración del arranque automático de los servicios}
\begin{lstlisting}
sudo systemctl enable docker.service
sudo systemctl enable containerd.service
\end{lstlisting}
    \item[§] \textbf{Comprobaciones}
    En caso de que no se haya presentado un problema en los pasos anteriores, se puede ejecutar el comando
\begin{lstlisting}
docker ps
\end{lstlisting}
\end{itemize}

% https://cursosdedesarrollo.com/2024/04/instalacion-de-docker-en-ubuntu-24-04-lts/